\documentclass[12pt,a4paper]{scrreprt}
\usepackage[utf8]{inputenc}
\usepackage[T1]{fontenc}
\usepackage{lmodern}
\usepackage{graphicx}
\usepackage{amsmath, amssymb}
\usepackage{natbib}
\usepackage{hyperref}
\bibliographystyle{apalike}

% ---------------------------------------------------------------
% Operator declarations (consistent naming throughout)
% ---------------------------------------------------------------
\DeclareMathOperator{\Mod}{Mod}
\DeclareMathOperator{\Med}{Med}
\DeclareMathOperator{\IOV}{IOV}
\DeclareMathOperator{\IQV}{IQV}
\DeclareMathOperator{\Bias}{Bias}

% ---------------------------------------------------------------
% TODO / open items:
%   - Bradley (2005) Fernleihe: Kapitel 5 (strong mixing conditions)
%   - Klenke, Bauer, Georgii (Maßtheorie), Brockwell, Shao (Math. Stats)
%   - Cramer-Wold justification for multivariate CLT
%   - Examples of categorical time series with exponential structure
%   - Oberwolfach reports
%   - Fokianos & Tjostheim (Paper)
%   - Section "Measures of Temporal Dependence" still empty
% ---------------------------------------------------------------

\begin{document}

% ================================================================
\chapter*{Categorical Time-Series Analysis}
% ================================================================

% ----------------------------------------------------------------
\section{General Setup}
% ----------------------------------------------------------------

\subsection{Notation}

Let $(X_t)_{t \in \mathbb{Z}}$ be a \textit{categorical} time series process, where each $X_t$
takes values in
\[
    S := \{s_0, \dots, s_m\}, \qquad m \in \mathbb{N}.
\]
$(X_t)$ may be \textit{nominal}, in which case the ordering of the elements of $S$ is arbitrary
and used solely for notational consistency, or \textit{ordinal}, in which case the ordering of
the elements of $S$ reflects an inherent order of the categories.

\subsection{Stationarity}

Two forms of stationarity are distinguished.
A process is called \emph{weakly stationary} if
\[
    \begin{aligned}
        \mathbb{E}(X_t)             &= \mu_X         &&\forall\, t,  \\
        \mathbb{E}(X_t^2)           &< \infty         &&\forall\, t,  \\
        \mathrm{Cov}(X_t, X_{t+h}) &= \gamma(h)      &&\forall\, t, h.
    \end{aligned}
\]
It is called \emph{strictly stationary} if all finite-dimensional distributions are invariant
under time shifts, i.e.,
\[
    F_{X_t, \ldots, X_{t+j}} = F_{X_{t+h}, \ldots, X_{t+h+j}}
    \qquad \forall\, t, j, h.
\]

% ----------------------------------------------------------------
\section{Marginal Distributions}
% ----------------------------------------------------------------

Regardless of whether one considers a simple categorical process with few categories or a
complex temporally evolving system, the analysis of categorical processes always rests on the
examination of the marginal probabilities. Because ordinal and nominal data encode different
types of information, their marginals are represented in different ways.

\paragraph{Ordinal data.}
Most analytical tools for ordinal processes exploit the natural ordering of the elements of $S$.
For this reason, the marginal distribution is represented by the cumulative distribution
function~(CDF), which accumulates probabilities in accordance with the ordering of the
categories:
\[
    \mathbf{f} = (f_0, \dots, f_{m-1})^{\top},
    \qquad
    f_i := P(X_t \leq s_i).
\]
It suffices to consider the cumulative probabilities of categories $s_0, \ldots, s_{m-1}$,
since $P(X_t \leq s_m) = 1$ by definition. Although one could equivalently estimate the PMF
and derive the CDF from it, the CDF representation is often more convenient because it is
directly used to compute location and dispersion measures.

\paragraph{Nominal data.}
Since nominal categories do not possess a meaningful order, statistical procedures cannot
exploit cumulative structure. The marginal distribution is therefore represented by the
probability mass function~(PMF):
\[
    \mathbf{p} = (p_0, \dots, p_{m})^{\top},
    \qquad
    p_i := P(X_t = s_i).
\]

% ----------------------------------------------------------------
\section{Population Measures}
% ----------------------------------------------------------------

\subsection{Location}

For a \textit{nominal} process any measure that relies on ordering is inappropriate. The
\emph{mode}, defined as the most probable category,
\[
    \Mod(X_t) = \arg\max_{s_i \in S}\, P(X_t = s_i),
\]
is therefore the natural choice, providing a meaningful summary by identifying the most
frequent category.

For an \textit{ordinal} process the natural ordering permits additional measures of central
tendency. While the mode remains informative, the \emph{median} respects the ordering and is
often more useful for describing the central location:
\[
    \Med(X_t) = \min\bigl\{\, s_i \in S : F(s_i) \geq \tfrac{1}{2} \,\bigr\}.
\]
In summary, the mode identifies the most typical category, whereas the median indicates the
midpoint on the ordered scale.

\subsection{Dispersion}

The concept of dispersion differs between ordinal and nominal processes. Two separate indices
are used: the \emph{Index of Ordinal Variation}~(IOV) for ordinal data and the \emph{Index of
Qualitative Variation}~(IQV) for nominal data. Both are scaled to $[0,1]$, where $0$ denotes
minimal dispersion and $1$ denotes maximal dispersion:
\[
    \IOV(\mathbf{f})
    = \frac{4}{m}\sum_{i=0}^{m-1} f_i (1 - f_i),
    \qquad
    \IQV(\mathbf{p})
    = \frac{m+1}{m}\!\left(1 - \sum_{i=0}^{m} p_i^2\right).
\]

\subsubsection{Minimal Dispersion}

Both measures attain the value $0$ if and only if the distribution is degenerate, i.e., all
probability mass is concentrated in a single category.

\paragraph{Ordinal case.}
A one-point distribution assigns all mass to a single category $s_j \in S$, corresponding to
a CDF vector of the form
\[
    \mathbf{f} \in
    \left\{
    \begin{pmatrix} 1 \\ 1 \\ \vdots \\ 1 \end{pmatrix},
    \begin{pmatrix} 0 \\ 1 \\ \vdots \\ 1 \end{pmatrix},
    \ldots,
    \begin{pmatrix} 0 \\ 0 \\ \vdots \\ 1 \end{pmatrix}
    \right\}
    =: \{\mathbf{c}_0,\ldots,\mathbf{c}_m\} \subset [0,1]^{m+1}.
\]
In each case every component satisfies $f_i \in \{0,1\}$, so $f_i(1-f_i) = 0$ for all $i$,
and hence $\IOV = 0$.

Conversely, $\IOV = 0$ implies $\sum_{i=0}^{m-1} f_i(1-f_i) = 0$. Since each term is
nonnegative, $f_i(1-f_i) = 0$ for all~$i$, which forces $f_i \in \{0,1\}$ for all~$i$.
Hence $\IOV = 0$ holds \emph{exactly} for one-point distributions.

\paragraph{Nominal case.}
A one-point distribution corresponds to a PMF vector of the form
\[
    \mathbf{p} \in
    \left\{
    \begin{pmatrix} 1 \\ 0 \\ \vdots \\ 0 \end{pmatrix},
    \begin{pmatrix} 0 \\ 1 \\ \vdots \\ 0 \end{pmatrix},
    \ldots,
    \begin{pmatrix} 0 \\ 0 \\ \vdots \\ 1 \end{pmatrix}
    \right\}
    =: \{\mathbf{e}_0,\ldots,\mathbf{e}_m\} \subset [0,1]^{m+1}.
\]
In each case $\sum_{i=0}^m p_i^2 = 1$, so $\IQV = 0$.

Conversely, $\IQV = 0$ implies $\sum_{i=0}^m p_i^2 = 1$. Since $0 \leq p_i \leq 1$, we have
$p_i^2 \leq p_i$ with equality if and only if $p_i \in \{0,1\}$. Combined with
$\sum_{i=0}^m p_i = 1$, exactly one component equals~$1$ and all others equal~$0$, i.e.,
$\mathbf{p} \in \{\mathbf{e}_0,\ldots,\mathbf{e}_m\}$.

\subsubsection{Maximal Dispersion}

\paragraph{Ordinal case.}
Dispersion is maximal when half the probability mass lies in the lowest category and half in
the highest, with all intermediate categories having probability zero. This yields
\[
    \mathbf{f} = \bigl(\tfrac{1}{2}, \tfrac{1}{2}, \ldots, \tfrac{1}{2}\bigr)^{\top},
\]
for which $f_i(1-f_i) = \tfrac{1}{4}$ for all $i$, giving $\IOV = 1$.

To confirm this is the unique maximiser, note that $g(f_i) := f_i(1-f_i)$ is strictly concave
with $g(0) = g(1) = 0$ and unique maximum at
\[
    \frac{\partial g}{\partial f_i} = 1 - 2f_i \overset{!}{=} 0
    \quad \Longrightarrow \quad
    f_i = \tfrac{1}{2}.
\]
Hence $\IOV(\mathbf{f}) = 1 \Leftrightarrow \mathbf{f} = \bigl(\tfrac{1}{2}, \ldots,
\tfrac{1}{2}\bigr)^{\top}$.

\paragraph{Nominal case.}
Dispersion is maximal when the probability mass is distributed uniformly across all $m+1$
categories:
\[
    p_i = \frac{1}{m+1}, \qquad i = 0,\ldots,m.
\]
Then $\sum_{i=0}^m p_i^2 = \tfrac{1}{m+1}$, so $\IQV = 1$.

The uniqueness of this maximiser follows from minimising $f(\mathbf{p}) := \sum_{i=0}^m p_i^2$
subject to $\sum_{i=0}^m p_i = 1$. The Lagrangian
\[
    \mathcal{L}(\mathbf{p},\lambda)
    := \sum_{i=0}^m p_i^2 - \lambda\!\left(\sum_{i=0}^m p_i - 1\right)
\]
yields first-order conditions $2p_j = \lambda$ for all $j$, so all components are equal. The
constraint then gives $p_j = \tfrac{1}{m+1}$ for all $j \in \{0,\ldots,m\}$.

\subsection{Skewness}

A simple measure for the skewness of an ordinal process is
\[
    \operatorname{skew}(\mathbf{f}) = \frac{2}{m} \sum_{i=0}^{m-1} f_i - 1.
\]

% ----------------------------------------------------------------
\section{Bivariate Lag Probabilities}
% ----------------------------------------------------------------

Temporal dependence in categorical time series can be quantified through bivariate
lag-$h$ probabilities, representing the joint probability that the process takes a specific
category at time~$t$ and another category at time~$t+h$.

\paragraph{Ordinal series (cumulative).}
\[
    f_{ij}(h) := P(X_t \leq s_i,\, X_{t+h} \leq s_j),
    \qquad i,j < m.
\]

\paragraph{Nominal series (joint).}
\[
    p_{ij}(h) := P(X_t = s_i,\, X_{t+h} = s_j),
    \qquad i,j \leq m.
\]

These quantities serve as input for dependence measures such as the ordinal and nominal
Cohen's~$\kappa$.

% ----------------------------------------------------------------
\section{Measures of Temporal Dependence}
% ----------------------------------------------------------------

% TODO: fill in this section (Cohen's kappa, ACF analogues, etc.)

% ================================================================
\chapter*{Sample Estimation}
% ================================================================

% ----------------------------------------------------------------
\section{Estimators under Complete Observation}
% ----------------------------------------------------------------

To estimate the marginal (cumulative) probabilities, the relative frequencies are expressed
through a binarisation of the original process $(X_t)$. For ordinal data, define the indicator
vector
\[
    \mathbf{Z}_t := (Z_{t,0}, \ldots, Z_{t,m-1})^{\top},
    \qquad
    Z_{t,i} := \mathbb{I}\{X_t \leq s_i\},
\]
and for nominal data analogously
\[
    \mathbf{Y}_t := (Y_{t,0}, \ldots, Y_{t,m})^{\top},
    \qquad
    Y_{t,i} := \mathbb{I}\{X_t = s_i\}.
\]
The marginal CDF and PMF vectors are then estimated by the corresponding sample means:
\[
    \hat{\mathbf{f}} := \frac{1}{n} \sum_{t=1}^n \mathbf{Z}_t,
    \qquad
    \hat{\mathbf{p}} := \frac{1}{n} \sum_{t=1}^n \mathbf{Y}_t.
\]

\subsection{Weak Dependence Assumptions}

To derive the asymptotic distribution of $\hat{\mathbf{f}}$ and $\hat{\mathbf{p}}$, a CLT for
dependent data is required. Following \citet{Bradley2005}, let
\[
    \mathcal{F}_J^L := \sigma(X_t : J \leq t \leq L)
\]
denote the $\sigma$-field generated by $\{X_t : J \leq t \leq L\}$. The $\alpha$-mixing
coefficients of $(X_t)$, given strict stationarity, are
\[
    \alpha_{X_t}(n)
    := \sup_{A \in \mathcal{F}_{-\infty}^0,\; B \in \mathcal{F}_{n}^{\infty}}
    \bigl| P(A \cap B) - P(A)P(B) \bigr|.
\]
The process is called $\alpha$-mixing if $\alpha(n) \to 0$ as $n \to \infty$. We additionally
assume \emph{exponentially decaying} mixing coefficients: there exist constants $C > 0$ and
$\rho \in (0,1)$ such that
\[
    \alpha(n) \leq C\rho^n \qquad \forall\, n \in \mathbb{N}.
\]

\subsection{Asymptotic Distribution of \texorpdfstring{$\hat{\mathbf{f}}$}{f-hat}}

The indicator mappings $g_i(X_t) := \mathbb{I}\{X_t \leq s_i\}$ and
$h_i(X_t) := \mathbb{I}\{X_t = s_i\}$ are measurable functions of $X_t$
(see \citealt{Klenke2005}, p.~34). Measurability preserves strict stationarity, so the
binarised processes $(g_i(X_t))$ and $(h_i(X_t))$ inherit strict stationarity from $(X_t)$.
Moreover,
\[
    \sigma(Z_{t,i} : J \leq t \leq L) \subset \sigma(X_t : J \leq t \leq L),
    \qquad J \leq L,
\]
so $\alpha_{Z_{t,i}}(n) \leq \alpha_{X_t}(n)$ for all $i$ and $n$, and hence $\mathbf{Z}_t$
and $\mathbf{Y}_t$ are also $\alpha$-mixing with exponentially decaying coefficients.

The applicable CLT is due to \citet{Ibragimov1962}, which requires strict stationarity,
$\alpha$-mixing, and the summability condition
\[
    \sum_{n=1}^{\infty} \alpha(n)^{\lambda/(2+\lambda)} < \infty
    \quad \text{for some } \lambda > 0.
\]
Under exponential decay this is satisfied:
\[
    \sum_{n=1}^{\infty} \alpha(n)^{\lambda/(2+\lambda)}
    \leq C^{\lambda/(2+\lambda)} \sum_{n=1}^{\infty} \rho^{\,n\lambda/(2+\lambda)}
    = C^{\lambda/(2+\lambda)} \frac{\rho^{\lambda/(2+\lambda)}}{1-\rho^{\lambda/(2+\lambda)}}
    < \infty.
\]
The moment condition $\mathbb{E}|Z_{t,i}|^{2+\lambda} < \infty$ is trivially satisfied since
$|Z_{t,i}|^{2+\lambda} \in \{0,1\}$.

\paragraph{Cramér–Wold argument.}
By the Cramér–Wold device, $(\mathbf{Z}_t)_t$ converges in distribution to a multivariate
normal if and only if every scalar projection
\[
    U_t := \mathbf{a}^{\top}\mathbf{Z}_t, \qquad \mathbf{a} \in \mathbb{R}^{m} \setminus \{\mathbf{0}\},
\]
converges to a univariate normal. Since $U_t = (\mathbf{a}^{\top} \circ g)(X_t)$ is a
measurable function of $X_t$, it inherits strict stationarity and $\alpha$-mixing. Moreover,
$|U_t| \leq M := \max_{j}\bigl|\sum_{i=0}^{j}a_i\bigr| < \infty$, so
$\mathbb{E}|U_t|^{2+\lambda} \leq M^{2+\lambda} < \infty$. Ibragimov's CLT is therefore
applicable to $U_t$, yielding
\[
    \sqrt{n}(\bar{U}_n - \mu_{U}) \xrightarrow{d} \mathcal{N}(0, \sigma_{\mathbf{a}}),
\]
where $\mu_U = \mathbb{E}[U_t]$ and
\[
    \sigma_{\mathbf{a}} = \sum_{i,j} a_i a_j \left[
        \mathrm{Cov}(Z_{t,i}, Z_{t,j})
        + \sum_{h=1}^{\infty}
        \Bigl(\mathrm{Cov}(Z_{t,i}, Z_{t+h,j}) + \mathrm{Cov}(Z_{t,j}, Z_{t+h,i})\Bigr)
    \right].
\]
By the Cramér–Wold device, the sample mean $\hat{\mathbf{f}} = \bar{\mathbf{Z}}_n$ satisfies
\[
    \sqrt{n}\bigl(\hat{\mathbf{f}} - \mathbf{f}\bigr)
    \xrightarrow{d} \mathcal{N}\bigl(\mathbf{0}, \boldsymbol{\Sigma}_Z\bigr),
\]
where the $(i,j)$-entry of $\boldsymbol{\Sigma}_Z$ is
\[
    \sigma_{Z;i,j}
    = \mathrm{Cov}(Z_{t,i}, Z_{t,j})
    + \sum_{h=1}^{\infty}
    \Bigl(\mathrm{Cov}(Z_{t,i}, Z_{t+h,j}) + \mathrm{Cov}(Z_{t,j}, Z_{t+h,i})\Bigr).
\]
Using $\mathrm{Cov}(Z_{t,i}, Z_{t,j}) = f_{\min\{i,j\}} - f_i f_j$ and
$\mathrm{Cov}(Z_{t,i}, Z_{t+h,j}) = f_{ij}(h) - f_i f_j$, this becomes
\[
    \sigma_{Z;i,j}
    = \bigl(f_{\min\{i,j\}} - f_i f_j\bigr)
    + \sum_{h=1}^{\infty} \bigl(f_{ij}(h) + f_{ji}(h) - 2f_i f_j\bigr).
\]

\subsection{Asymptotic Distribution of Sample Measures}

Given the asymptotic normality of $\hat{\mathbf{f}}$, the delta method can be applied to
derive the asymptotic properties of sample measures that are smooth functions of $\hat{\mathbf{f}}$.

\paragraph{IOV.}
The Jacobian of $\IOV(\mathbf{f})$ with respect to $\mathbf{f}$ is
\[
    \mathbf{j}_{\IOV}
    = \frac{4}{m}\bigl(1-2f_0,\, \ldots,\, 1-2f_{m-1}\bigr),
\]
giving asymptotic variance
\[
    \mathrm{Var}\bigl(\widehat{\IOV}\bigr)
    \overset{\mathrm{asy.}}{=}
    \mathbf{j}_{\IOV}\,\boldsymbol{\Sigma}_Z\,\mathbf{j}_{\IOV}^{\top}
    = \frac{16}{m^2} \sum_{i,j=0}^{m-1}(1-2f_i)(1-2f_j)\,\sigma_{Z;i,j}.
\]
Since $\IOV$ is nonlinear, the delta method only characterises the asymptotic distribution; a
finite-sample bias of order $n^{-1}$ is present. The Hessian is
\[
    \mathbf{H}_{\IOV}(\mathbf{f})
    = \left(\frac{\partial^2 \IOV}{\partial f_i \partial f_j}\right)_{i,j}
    = -\frac{8}{m}\,\mathbf{I}_m,
\]
and a second-order Taylor expansion yields
\[
    \Bias\bigl(\widehat{\IOV}\bigr)
    = \frac{1}{2}\operatorname{tr}\bigl(\mathbf{H}_{\IOV}\,\mathrm{Cov}(\hat{\mathbf{f}})\bigr)
    = -\frac{4}{m}\sum_{i=0}^{m-1}\sigma_{Z;i,i}.
\]
The bias is strictly negative and depends only on the marginal variances $\sigma_{Z;i,i}$.

\paragraph{Skewness.}
The Jacobian of $\operatorname{skew}(\mathbf{f})$ is
\[
    \mathbf{j}_{\mathrm{skew}}
    = \left(\frac{2}{m},\, \ldots,\, \frac{2}{m}\right),
\]
giving
\[
    \mathrm{Var}\bigl(\widehat{\operatorname{skew}}\bigr)
    \overset{\mathrm{asy.}}{=}
    \frac{4}{m^2}\sum_{i,j=0}^{m-1}\sigma_{Z;ij},
    \qquad
    \mathbb{E}\bigl[\widehat{\operatorname{skew}}\bigr]
    \overset{\mathrm{asy.}}{=} \operatorname{skew}(\mathbf{f}).
\]
Since skewness is linear in $\mathbf{f}$, the Hessian vanishes and the estimator is exactly
unbiased for all~$n$.

% ----------------------------------------------------------------
\section{Estimation under Missingness}
% ----------------------------------------------------------------

\subsection{Amplitude Modulation Model}

Missingness is modelled via \emph{amplitude modulation}. A binary stochastic process
$(O_t)_{t \in \mathbb{Z}}$ runs concurrently with $(X_t)$ and modulates observations:
\[
    O_t \in \{0,1\}, \qquad O_t = 1 \;\Leftrightarrow\; X_t \text{ is observed}.
\]
The amplitude-modulating process $(O_t)$ may be deterministic ($O_t = o_t$, known) or
stochastic with $\pi_t := P(O_t = 1)$, possibly depending on the history of $(X_t)$, $(O_t)$,
or external factors.

The observable process is $(O_t\mathbf{Z}_t)$, and the \emph{naive estimator} of $\mathbf{f}$ is
\[
    \hat{\mathfrak{f}} := \frac{1}{n}\sum_{t=1}^{n} O_t\mathbf{Z}_t.
\]
Its expectation, by the law of iterated expectations, satisfies
\[
    \mathbb{E}[\hat{\mathfrak{f}}]
    = \frac{1}{n}\sum_{t=1}^{n} \mathbb{E}\bigl[\mathbf{Z}_t\,\mathbb{E}[O_t \mid \mathbf{Z}_t]\bigr].
\]
Under the assumption that $(O_t)$ and $(X_t)$ are independent, this simplifies to
\[
    \mathbb{E}[\hat{\mathfrak{f}}]
    = \left(\frac{1}{n}\sum_{t=1}^{n}\mathbb{E}[O_t]\right)\mathbf{f}.
\]
The naive estimator is therefore biased towards zero by the factor
$\bar{\pi} := \tfrac{1}{n}\sum_{t=1}^n \mathbb{E}[O_t]$. This motivates the \emph{corrected
estimator}
\[
    \hat{\mathbf{f}}^{*} := \frac{\hat{\mathfrak{f}}}{\bar{O}}
    = \frac{\tfrac{1}{n}\sum_{t=1}^{n} O_t\mathbf{Z}_t}{\tfrac{1}{n}\sum_{t=1}^{n}O_t}.
\]

If $(O_t)$ is deterministic then $\bar{O} = \bar{\pi}$ and $\hat{\mathbf{f}}^{*}$ is exactly
unbiased. If $(O_t)$ is stochastic then $\bar{O} \neq \bar{\pi}$ in general, and
$\hat{\mathbf{f}}^{*}$ carries a finite-sample bias studied below.

\subsection{Asymptotic Distribution of \texorpdfstring{$\hat{\mathbf{f}}^{*}$}{f*-hat}}

Assume $(O_t)$ and $(X_t)$ are independent, stationary, and $\alpha$-mixing with exponentially
decaying coefficients. Write $\mathbb{E}[O_t] = \pi$, $\gamma_O(h) := \mathrm{Cov}(O_h, O_0)$,
and $\pi(h) := \mathbb{E}[O_h O_0] = \gamma_O(h) + \pi^2$.

Introduce the extended binary vector
\[
    \mathbf{Z}_t^{*} := O_t(1, \mathbf{Z}_t^{\top})^{\top}
    = (Z_{t,-1}^{*}, Z_{t,0}^{*}, \ldots, Z_{t,m-1}^{*})^{\top},
\]
with expectation
\[
    \mathbb{E}[\mathbf{Z}_t^{*}]
    = \pi(1, \mathbf{f}^{\top})^{\top}
    =: \boldsymbol{\mu}_{\mathbf{Z}^{*}}.
\]
By the same Cramér–Wold argument as before, the sample mean
$\bar{\mathbf{Z}}_n^{*} := \tfrac{1}{n}\sum_{t=1}^n \mathbf{Z}_t^{*}$ satisfies
\[
    \sqrt{n}\bigl(\bar{\mathbf{Z}}_n^{*} - \boldsymbol{\mu}_{\mathbf{Z}^{*}}\bigr)
    \xrightarrow{d} \mathcal{N}\bigl(\mathbf{0}, \boldsymbol{\Sigma}_{\mathbf{Z}^{*}}\bigr),
\]
where $\boldsymbol{\Sigma}_{\mathbf{Z}^{*}} = \mathrm{Var}(\mathbf{Z}_t^{*}) + 2\sum_{h=1}^{\infty}\mathrm{Cov}(\mathbf{Z}_t^{*}, \mathbf{Z}_{t+h}^{*})$. The individual entries are:

\medskip
\noindent\textbf{(i)} $i = j = -1$:
\[
    \sigma_{\mathbf{Z}^{*};-1,-1}
    = \pi(1-\pi) + 2\sum_{h=1}^{\infty}\gamma_O(h).
\]

\noindent\textbf{(ii)} $j \geq 0,\; i = -1$ (by independence of $(O_t)$ and $(X_t)$):
\[
    \sigma_{\mathbf{Z}^{*};-1,j}
    = f_j\,\sigma_{\mathbf{Z}^{*};-1,-1}.
\]

\noindent\textbf{(iii)} $i, j \geq 0$:
\[
    \sigma_{\mathbf{Z}^{*};i,j}
    = f_i f_j\,\sigma_{\mathbf{Z}^{*};-1,-1}
    + \pi\bigl(f_{\min\{i,j\}} - f_i f_j\bigr)
    + \sum_{h=1}^{\infty}\pi(h)\bigl(f_{ij}(h) + f_{ji}(h) - 2f_i f_j\bigr).
\]

\paragraph{Delta method for $\hat{\mathbf{f}}^{*}$.}
The corrected estimator is $\hat{\mathbf{f}}^{*} = \mathbf{g}(\bar{\mathbf{Z}}_n^{*})$ where
$g_j(x_{-1}, x_0, \ldots, x_{m-1}) = x_j / x_{-1}$. The Jacobian of $\mathbf{g}$ evaluated
at $\boldsymbol{\mu}_{\mathbf{Z}^{*}}$ is
\[
    \mathbf{G} =
    \begin{pmatrix}
        -f_0/\pi    & 1/\pi  & 0      & \cdots & 0      \\
        -f_1/\pi    & 0      & 1/\pi  & \ddots & \vdots \\
        \vdots      & \vdots & \ddots & \ddots & 0      \\
        -f_{m-1}/\pi& 0      & \cdots & 0      & 1/\pi
    \end{pmatrix}
    \in \mathbb{R}^{m \times (m+1)}.
\]
By the multivariate delta method,
\[
    \sqrt{n}\bigl(\hat{\mathbf{f}}^{*} - \mathbf{f}\bigr)
    \xrightarrow{d} \mathcal{N}\bigl(\mathbf{0}, \boldsymbol{\Sigma}_{\hat{\mathbf{f}}^{*}}\bigr),
    \qquad
    \boldsymbol{\Sigma}_{\hat{\mathbf{f}}^{*}} = \mathbf{G}\,\boldsymbol{\Sigma}_{\mathbf{Z}^{*}}\,\mathbf{G}^{\top}.
\]
Carrying out the matrix multiplication, the $(i,j)$-entry simplifies to
\[
    \sigma_{\hat{\mathbf{f}}^{*};i,j}
    = \frac{1}{\pi}\bigl(f_{\min\{i,j\}} - f_i f_j\bigr)
    + \frac{1}{\pi^2}\sum_{h=1}^{\infty}\pi(h)
    \bigl(f_{ij}(h) + f_{ji}(h) - 2f_i f_j\bigr).
\]

\paragraph{Bias of $\hat{\mathbf{f}}^{*}$.}
The Hessian of $g_j$ at $\boldsymbol{\mu}_{\mathbf{Z}^{*}}$ has only two nonzero entries:
\[
    H_{g_j;-1,-1} = \frac{2f_j}{\pi^2}, \qquad
    H_{g_j;-1,j} = H_{g_j;j,-1} = -\frac{1}{\pi^2}.
\]
The second-order bias is then
\[
    \Bias(\hat{f}_j^{*})
    = \frac{1}{2}\operatorname{tr}\bigl(\mathbf{H}_{g_j}\,\boldsymbol{\Sigma}_{\mathbf{Z}^{*}}\bigr)
    = \frac{f_j}{\pi^2}\,\sigma_{\mathbf{Z}^{*};-1,-1}
    - \frac{1}{\pi^2}\cdot f_j\,\sigma_{\mathbf{Z}^{*};-1,-1}
    = 0,
\]
using $\sigma_{\mathbf{Z}^{*};-1,j} = f_j\,\sigma_{\mathbf{Z}^{*};-1,-1}$ from entry~(ii) above.
Hence $\hat{\mathbf{f}}^{*}$ is asymptotically unbiased to second order.

\paragraph{Special cases.}
Three illuminating special cases are worth noting:

\begin{enumerate}
    \item \textit{Full observation} ($\pi = 1$): $\sigma_{\hat{\mathbf{f}}^{*};i,j} = \sigma_{Z;i,j}$,
    recovering the complete-data asymptotic variance.

    \item \textit{i.i.d.\ missingness} ($\pi(h) = \pi^2$ for all $h \geq 1$):
    \[
        \sigma_{\hat{\mathbf{f}}^{*};i,j}
        = \frac{1}{\pi}\bigl(f_{\min\{i,j\}} - f_i f_j\bigr)
        + \sum_{h=1}^{\infty}\bigl(f_{ij}(h) + f_{ji}(h) - 2f_i f_j\bigr).
    \]

    \item \textit{i.i.d.\ process} ($f_{ij}(h) = f_i f_j$ for all $h \geq 1$):
    \[
        \sigma_{\hat{\mathbf{f}}^{*};i,j}
        = \frac{1}{\pi}\bigl(f_{\min\{i,j\}} - f_i f_j\bigr).
    \]
\end{enumerate}

\subsection{Asymptotic Properties of Marginal Measures under Missingness}

The delta method can be applied to $\hat{\mathbf{f}}^{*}$ in the same way as to $\hat{\mathbf{f}}$,
yielding the asymptotic distribution of measures such as $\widehat{\IOV}^{*}$. The corresponding
Jacobians and Hessians are identical to those computed in Section~3, with $\boldsymbol{\Sigma}_Z$
replaced by $\boldsymbol{\Sigma}_{\hat{\mathbf{f}}^{*}}$.

\bibliography{references}

\end{document}
